\let\R\relax
\newcommand{\R}{\mathbb{R}}
\let\C\relax
\newcommand{\C}{\mathbb{C}}
\let\N\relax
\newcommand{\N}{\mathbb{N}}
\let\Z\relax
\newcommand{\Z}{\mathbb{Z}}
\let\Q\relax
\newcommand{\Q}{\mathbb{Q}}
\let\F\relax
\newcommand{\F}{\mathbb{F}}
\let\SO\relax
\newcommand{\SO}{\operatorname{SO}}
\let\SL\relax
\newcommand{\SL}{\operatorname{SL}}
\let\GL\relax
\newcommand{\GL}{\operatorname{GL}}
\let\St\relax
\newcommand{\St}{\operatorname{St}}
\let\GV\relax
\newcommand{\GV}{\operatorname{GV}}
\let\E\relax
\newcommand{\E}{\operatorname{E}}
\let\Gal\relax
\newcommand{\Gal}{\operatorname{Gal}}

\chapter{Andrew Wiles (1995/96)}
\index{Wiles, Andrew}

\begin{quote}
\it ``Wiles's proof of Fermat's Last Theorem is a landmark of intellectual history. It is not only the resolution of a 350-year-old problem but also a testament to the power of the interconnectedness of mathematics, bringing together modular forms, elliptic curves, Galois representations, and the Langlands program in a synthesis that few could have imagined possible. His work has fundamentally changed the landscape of number theory.'' --- Wolf Prize Citation
\end{quote}

Andrew John Wiles (born 11 April 1953) is one of the most celebrated mathematicians of the modern era. In 1994, he proved Fermat's Last Theorem, a problem that had remained unsolved for over 350 years and was arguably the most famous open problem in mathematics. For this achievement, he received the 1995/96 Wolf Prize in Mathematics, shared with Robert Langlands. He was later awarded the 2016 Abel Prize.

Wiles's proof is a monumental synthesis of 20th-century number theory. It draws on the theory of elliptic curves, modular forms, Galois representations, Iwasawa theory, and the Langlands program. The proof, completed with the assistance of Richard Taylor, established the modularity of semistable elliptic curves over \(\Q\), a special case of the Shimura--Taniyama--Weil conjecture. The resolution of Fermat's Last Theorem was not merely the solution of an ancient puzzle; it was a triumph of the deep structural vision that modern arithmetic geometry had developed over the preceding decades.

\section{Early Life and Career}

Andrew Wiles was born on 11 April 1953 in Cambridge, England, to Maurice Wiles (a professor of divinity at Oxford) and Patricia Wiles. He grew up in Cambridge and developed an interest in mathematics at an early age. He has recounted being fascinated by Fermat's Last Theorem after reading about it in a library book at age ten, resolving to be the first to prove it.

Wiles attended King's College School, Canterbury, and then The Leys School, Cambridge. In 1971, he entered Merton College, Oxford, where he read mathematics. He graduated with a first-class degree in 1974. He pursued graduate studies at Clare College, Cambridge, under the supervision of John Coates. His doctoral work was on the arithmetic of elliptic curves and Iwasawa theory, subjects that would prove essential to his later work on Fermat's Last Theorem.

Wiles completed his PhD in 1979 with a thesis titled \textit{Reciprocity Laws and the Conjecture of Birch and Swinnerton-Dyer}. He then held a Junior Research Fellowship at Clare College (1977--1980) before moving to the United States. He was a visiting scholar at the Institute for Advanced Study in Princeton (1980) and then joined the faculty at Harvard University (1981--1982). In 1982, he became a professor at Princeton University, where he remained for most of his career, apart from a period at the Institute for Advanced Study (1994--1996) and later at Oxford.

Wiles's early work focused on the Birch and Swinnerton-Dyer conjecture, complex multiplication, and Iwasawa theory. He proved important results on the Goldfeld--Szpiro conjecture and the Mazur--Wiles theorem on the main conjecture of Iwasawa theory for cyclotomic fields (with Barry Mazur). These deep results in arithmetic geometry laid the foundation for his later assault on Fermat's Last Theorem.

\section{Fermat's Last Theorem: A Historical Overview}

Fermat's Last Theorem originated from a marginal note made by Pierre de Fermat around 1637 in his copy of Diophantus's \textit{Arithmetica}. Fermat wrote that he had discovered a truly marvelous proof that the equation

\[
x^n + y^n = z^n
\]

has no positive integer solutions for \(n \geq 3\), but that the margin was too small to contain it. This tantalizing claim became the most famous unsolved problem in mathematics.

\subsection{Early Progress}

For over 350 years, mathematicians attempted to prove Fermat's statement. Significant progress was made on special cases:

\begin{itemize}
    \item \textbf{n = 4:} Fermat himself proved the case \(n = 4\) using the method of infinite descent, showing that \(x^4 + y^4 = z^4\) has no nontrivial solutions.
    \item \textbf{n = 3:} Leonhard Euler (1770) proved the case \(n = 3\) using the method of infinite descent applied to numbers of the form \(a + b\sqrt{-3}\).
    \item \textbf{n = 5:} Dirichlet and Legendre (1825) independently proved the case \(n = 5\).
    \item \textbf{n = 7:} Lam\'e (1839) proved the case \(n = 7\).
    \item \textbf{Sophie Germain:} In the early 19th century, Germain proved that if \(p\) is an odd prime such that \(2p+1\) is also prime, then there are no solutions to \(x^p + y^p = z^p\) with \(xyz\) not divisible by \(p\).
    \item \textbf{Kummer:} Ernst Kummer (1850) developed the theory of ideal numbers to prove FLT for all regular primes, establishing a deep connection between FLT and algebraic number theory.
\end{itemize}

Despite these advances, a general proof remained elusive. The problem resisted the efforts of the greatest mathematicians for centuries.

\subsection{The Modern Era}

The modern approach to Fermat's Last Theorem emerged from a series of profound insights in the 1960s--1980s that connected the problem to the theory of elliptic curves and modular forms. The key steps were:

\begin{enumerate}
    \item The construction by Gerhard Frey (1985) of an elliptic curve associated to a putative solution of Fermat's equation.
    \item The proof by Ken Ribet (1986) that Frey's curve, if it existed, could not be modular, thereby showing that the Shimura--Taniyama--Weil conjecture (stating that every elliptic curve over \(\Q\) is modular) implies Fermat's Last Theorem.
    \item Andrew Wiles's proof (1994) of the modularity of semistable elliptic curves, a sufficient case of the Shimura--Taniyama--Weil conjecture to establish Fermat's Last Theorem.
\end{enumerate}

\section{The Frey Curve}

The crucial insight connecting Fermat's Last Theorem to modular forms came from Gerhard Frey in 1985. Frey proposed that if a nontrivial solution to Fermat's equation existed, then a certain associated elliptic curve would have properties so pathological that it could not exist.

\begin{definition}[Frey Curve]
Let \(p \geq 5\) be a prime, and suppose \(a, b, c\) are positive integers with
\[
a^p + b^p = c^p.
\]
The \textbf{Frey curve} (or Frey--Hellegouarch curve) associated to this triple is the elliptic curve over \(\Q\) defined by the Weierstrass equation
\[
E_{a,b}: \quad y^2 = x(x - a^p)(x + b^p).
\]
\end{definition}

The discriminant of this curve is
\[
\Delta = (abc)^{2p} / 256,
\]
and its conductor is divisible only by the primes dividing \(abc\). The key properties of the Frey curve are:

\begin{itemize}
    \item The curve \(E_{a,b}\) is semistable (its bad reduction is multiplicative rather than additive) at all primes.
    \item The discriminant \(\Delta\) is a perfect \(p\)-th power up to a small factor.
    \item The Galois representation on the \(p\)-torsion points of \(E_{a,b}\) has very restricted ramification.
\end{itemize}

Frey argued that these properties are incompatible with the curve being modular. Specifically, if \(E_{a,b}\) were modular, associated to a modular form \(f\) of level \(N\), then level-lowering theorems would force the existence of a modular form of much lower level --- ultimately level 2 --- which Ribet proved impossible.

\section{Ribet's Theorem: Level Lowering}

Ken Ribet, building on ideas of Barry Mazur and work of Serre, proved the critical theorem that completed the connection between Fermat's Last Theorem and modularity.

\begin{theorem}[Ribet's Theorem, 1986]
Let \(E\) be an elliptic curve over \(\Q\) with conductor \(N\). Suppose that the mod \(p\) Galois representation
\[
\overline{\rho}_{E,p}: \Gal(\overline{\Q}/\Q) \to \GL_2(\F_p)
\]
is modular of level \(N\) (arising from a modular form of level \(N\)). If \(p\) is a prime such that the representation is irreducible and certain technical conditions hold, then \(\overline{\rho}_{E,p}\) is also modular of level \(N'\), where \(N'\) divides \(N/p\) for every prime \(\ell \parallel N\) at which the representation is finite at \(\ell\).

In particular, for the Frey curve \(E_{a,b}\) associated to a putative Fermat solution, Ribet showed that the level of the associated modular form could be reduced to 2.
\end{theorem}

Ribet's theorem uses the theory of congruences between modular forms, developed in the work of Mazur, Hida, and others. The essential idea is this: if a Galois representation \(\overline{\rho}\) arises from a modular form \(f\) of level \(N\), then for each prime \(\ell \mid N\) where the representation is unramified, one can find a congruent modular form of level \(N/\ell\). By iterating this process, the level can be stripped of all primes dividing \(abc\), until only primes dividing 2 remain.

The final step is a theorem of Mazur that there are no non-trivial modular forms of level 2 that could give rise to the representation \(\overline{\rho}_{E,p}\). Therefore, if \(E_{a,b}\) were modular, we would reach a contradiction. Hence:

\begin{corollary}
If every semistable elliptic curve over \(\Q\) is modular, then Fermat's Last Theorem holds.
\end{corollary}

This corollary is what motivated Wiles to prove the modularity of semistable elliptic curves. The problem had been reduced from a 350-year-old number theory puzzle to a contemporary question about the modularity of elliptic curves.

\section{The Shimura--Taniyama--Weil Conjecture}

The conjecture that every elliptic curve over \(\Q\) is modular originated in the work of Yutaka Taniyama, Goro Shimura, and Andr\'e Weil in the 1950s and 1960s.

\subsection{Elliptic Curves over \(\Q\)}

An elliptic curve over \(\Q\) is a smooth projective curve of genus 1 with a distinguished rational point, typically given by a Weierstrass equation
\[
E: \quad y^2 = x^3 + ax + b,
\]
with \(a, b \in \Q\) and discriminant \(\Delta_E = -16(4a^3 + 27b^2) \neq 0\). Every elliptic curve has a conductor \(N_E\), an integer encoding its bad reduction behavior.

The set of rational points \(E(\Q)\) forms a finitely generated abelian group (Mordell's theorem). The Birch and Swinnerton-Dyer conjecture relates the rank of this group to the behavior of the \(L\)-function of \(E\) at \(s = 1\).

\subsection{Modular Forms}

A modular form of weight \(k\) and level \(N\) is a holomorphic function \(f: \mathcal{H} \to \C\) on the upper half-plane that transforms under the action of the congruence subgroup \(\Gamma_0(N)\) in a specific way:
\[
f\left(\frac{az + b}{cz + d}\right) = (cz + d)^k f(z),
\]
for all \(\bigl(\begin{smallmatrix}a&b\\c&d\end{smallmatrix}\bigr) \in \Gamma_0(N)\). It has a Fourier expansion at infinity:
\[
f(z) = \sum_{n=0}^\infty a_n q^n, \qquad q = e^{2\pi i z}.
\]

A \textbf{newform} is a modular form that is an eigenvector for all Hecke operators and is not induced from a lower level. To a newform \(f\) of weight 2, one associates an elliptic curve \(E_f\) over \(\Q\) via the Eichler--Shimura construction, such that
\[
L(E_f, s) = \sum_{n=1}^\infty \frac{a_n}{n^s}.
\]

\begin{condition}[Shimura--Taniyama--Weil Conjecture, circa 1960]
Every elliptic curve over \(\Q\) is \textbf{modular}: for every elliptic curve \(E/\Q\) of conductor \(N\), there exists a newform \(f\) of weight 2 and level \(N\) such that \(L(E,s) = L(f,s)\). Equivalently, there is a surjective morphism
\[
X_0(N) \to E,
\]
where \(X_0(N)\) is the modular curve parameterizing elliptic curves with level \(N\) structure.
\end{condition}

This conjecture was considered extremely deep and far beyond the reach of available techniques. It became a central problem in number theory, with connections to the Langlands program, the theory of automorphic forms, and arithmetic geometry.

\section{Wiles's Strategy: Modularity Lifting}

Wiles's fundamental insight was to approach the modularity conjecture through the deformation theory of Galois representations. The strategy is to show that if the mod \(p\) representation \(\overline{\rho}_{E,p}\) is modular, then the full \(p\)-adic representation \(\rho_{E,p}\) is modular. This is a \textbf{modularity lifting} theorem.

\subsection{Galois Representations Associated to Elliptic Curves}

Let \(E\) be an elliptic curve over \(\Q\). For any prime \(p\), the \(p\)-torsion points \(E[p] \cong (\Z/p\Z)^2\) form a Galois module, giving a representation
\[
\overline{\rho}_{E,p}: \Gal(\overline{\Q}/\Q) \to \GL_2(\F_p).
\]

Similarly, the Tate module \(T_p(E) = \varprojlim E[p^n] \cong \Z_p^2\) gives a \(p\)-adic representation
\[
\rho_{E,p}: \Gal(\overline{\Q}/\Q) \to \GL_2(\Z_p).
\]

If \(E\) is modular, associated to a newform \(f\), then there is a corresponding Galois representation \(\rho_f\) constructed by Deligne. The modularity of \(E\) is equivalent to the statement that \(\rho_{E,p} \cong \rho_f\) for some \(f\).

\subsection{Deformation Rings}

Wiles introduced the concept of \textbf{universal deformation rings} for Galois representations. Given a residual representation
\[
\overline{\rho}: \Gal(\overline{\Q}/\Q) \to \GL_2(k),
\]
where \(k\) is a finite field, one can consider all liftings of \(\overline{\rho}\) to complete Noetherian local \(\Z_p\)-algebras with residue field \(k\). Under suitable conditions, there exists a \textbf{universal deformation}
\[
\rho^{\text{univ}}: \Gal(\overline{\Q}/\Q) \to \GL_2(R),
\]
where \(R\) is the universal deformation ring, and every lifting of \(\overline{\rho}\) factors uniquely through \(\rho^{\text{univ}}\).

The deformation ring \(R\) is characterized by the property that it represents the deformation functor on the category of local Artinian \(\Z_p\)-algebras. Its structure encodes all possible liftings of the residual representation.

\subsection{Hecke Algebras}

Parallel to the deformation ring, Wiles considered the \textbf{Hecke algebra} \(\mathbb{T}\) acting on the space of modular forms of a given level. The Hecke algebra is a commutative \(\Z_p\)-algebra generated by Hecke operators \(T_\ell\). The Galois representations associated to modular forms give a map
\[
\mathbb{T} \to R,
\]
and the modularity lifting problem reduces to showing that this map is an isomorphism.

Wiles's modularity lifting theorem can be stated as follows:

\begin{theorem}[Wiles's Modularity Lifting Theorem, 1994]
Let \(p\) be an odd prime. Let \(\rho: \Gal(\overline{\Q}/\Q) \to \GL_2(\overline{\Z}_p)\) be a continuous irreducible representation that is unramified outside a finite set of primes and satisfies the following conditions:
\begin{enumerate}
    \item \(\rho\) is ordinary or flat at \(p\) in the sense of Fontaine--Mazur.
    \item The residual representation \(\overline{\rho}\) is absolutely irreducible and modular.
    \item \(\overline{\rho}\) is not induced from a character of a quadratic extension (i.e., it is not dihedral).
    \item \(\overline{\rho}\) has big image in the sense that its image contains \(\SL_2(\F_p)\) for odd \(p\) (or more generally, the representation is \textit{modular} with a mild condition).
\end{enumerate}
Then \(\rho\) is modular. In other words, if the mod \(p\) reduction of a \(p\)-adic Galois representation is modular, then the full \(p\)-adic representation is modular.
\end{theorem}

This theorem is the heart of Wiles's proof. To apply it to the Frey curve \(E = E_{a,b}\), one needs to know that its mod \(p\) representation \(\overline{\rho}_{E,p}\) is modular. This is not known a priori; however, Wiles chose a different prime and a different curve. He used the fact that for the semistable curves arising from the Fermat equation, one can work with a particular prime \(p\) for which the modularity of \(\overline{\rho}\) can be established.

\section{The Taylor--Wiles Method}

The proof of the modularity lifting theorem required a novel technique, developed by Wiles in collaboration with Richard Taylor, now known as the \textbf{Taylor--Wiles method} or the method of \textbf{deformation rings and Hecke algebras}.

\subsection{The Numerical Criterion}

The key idea is to compare the sizes of the deformation ring \(R\) and the Hecke algebra \(\mathbb{T}\). Wiles proved that under certain conditions, there is a surjection \(R \to \mathbb{T}\). To prove it is an isomorphism, one must show that \(R\) and \(\mathbb{T}\) have the same size, measured by their \textbf{canonical ideal} and \textbf{congruence ideal}.

More precisely, Wiles introduced a numerical criterion: if there is a surjection \(R \to \mathbb{T}\) of complete intersection rings, and if the lengths of certain quotients satisfy an equality, then the map is an isomorphism. The condition involves the cotangent space of the deformation ring and the congruence ideal of the Hecke algebra.

\subsection{Patching Over Siegel Units}

To establish the numerical criterion, Taylor and Wiles introduced a brilliant technique known as \textbf{patching} or the \textbf{Taylor--Wiles patching argument}. The idea is to consider not just one deformation problem but a sequence of auxiliary deformation problems with additional level structure.

Specifically, one introduces auxiliary primes \(\ell_1, \dots, \ell_r\) and considers modular forms of higher level \(N \ell_1 \cdots \ell_r\). For each such set of primes, there is a deformation ring \(R_n\) and a Hecke algebra \(\mathbb{T}_n\). As the number of primes increases (in a carefully controlled way), the corresponding rings form a projective system, and at the limit one obtains an isomorphism
\[
R_\infty \cong \mathbb{T}_\infty,
\]
where both are power series rings over \(\Z_p\) in several variables.

The patching argument proceeds as follows:

\begin{enumerate}
    \item Choose a finite set of primes \(\mathcal{Q} = \{\ell_1, \dots, \ell_r\}\) with the property that each \(\ell_i \equiv 1 \pmod{p^N}\) and that the deformation problem with ramification allowed at these primes is unobstructed in a precise sense.
    \item For each subset of \(\mathcal{Q}\), one defines a deformation ring \(R_{\mathcal{Q}}\) and a Hecke algebra \(\mathbb{T}_{\mathcal{Q}}\).
    \item The map \(R_{\mathcal{Q}} \to \mathbb{T}_{\mathcal{Q}}\) is a surjection of finite flat \(\Z_p\)-algebras.
    \item Taking the limit over larger and larger sets \(\mathcal{Q}\) yields an isomorphism \(R_\infty \cong \mathbb{T}_\infty\).
    \item By specializing back to the original deformation problem, one deduces that the original map \(R \to \mathbb{T}\) is an isomorphism.
\end{enumerate}

This patching argument is one of the most innovative techniques in modern number theory. It has since been applied to many other modularity problems, including the full Shimura--Taniyama--Weil conjecture (proved by Breuil, Conrad, Diamond, and Taylor in 2001) and many cases of the Langlands correspondence for \(\GL_2\) over number fields.

\section{The 1993 Announcement and the Gap}

By the spring of 1993, after seven years of intense, secret work, Wiles believed he had a complete proof. He announced his results at a series of three lectures at the Isaac Newton Institute in Cambridge, England, from June 21--23, 1993. The audience included many of the world's leading number theorists.

The announcement created a sensation. News of Wiles's proof spread around the world, making headlines in newspapers and magazines. It seemed that the most famous problem in mathematics had finally been solved.

However, during the rigorous refereeing process, a gap was discovered in the portion of the proof dealing with the study of the Hecke algebra at primes dividing the conductor. Specifically, the argument that the deformation ring \(R\) was a complete intersection (which was needed to apply the numerical criterion) turned out to have a flaw.

The gap was in the construction of the Euler system used to bound the size of the Hecke algebra. Wiles had attempted to bypass certain difficulties in the Kolyvagin--Flach approach, but the argument was incomplete. The error was identified by Nick Katz, who was helping Wiles review the manuscript.

Wiles spent the next several months attempting to fix the proof. By December 1993, he had made progress but had not yet resolved the issue. He publicly acknowledged the gap in a preprint circulated that month, stating that the situation was still under investigation.

\section{The 1994 Breakthrough}

The breakthrough came in September 1994, when Wiles realized that the method that had failed for the ring of Hecke operators at the full level could be replaced by a different approach. The key insight was to return to the original ideas of the Taylor--Wiles method and to combine them with the theory of \textbf{universal deformations} in a more refined way.

Wiles invited Richard Taylor, a former PhD student of his (then at Cambridge), to help develop the revised argument. In a crucial flash of insight, Wiles realized that the issues with the Euler system could be circumvented by extending the Taylor--Wiles patching argument to handle the case where the Hecke algebra was not a complete intersection a priori.

The revised proof was completed in October 1994. Two papers were submitted to the \textit{Annals of Mathematics}:

\begin{enumerate}
    \item A.\ Wiles, \textit{Modular Elliptic Curves and Fermat's Last Theorem} (109 pages).
    \item R.\ Taylor and A.\ Wiles, \textit{Ring-Theoretic Properties of Certain Hecke Algebras} (47 pages).
\end{enumerate}

These papers were published in the Annals in 1995. The proof was scrutinized by the mathematical community and universally accepted as correct.

\section{The Modularity Theorem (Semistable Case)}

The main result proved by Wiles is the modularity of all semistable elliptic curves over \(\Q\).

\begin{theorem}[Modularity Theorem --- Semistable Case, Wiles 1995]
Every semistable elliptic curve \(E\) over \(\Q\) is modular. That is, for every elliptic curve \(E/\Q\) with semistable reduction (multiplicative reduction at all primes of bad reduction), there exists a cuspidal newform \(f\) of weight 2 and level \(N\) equal to the conductor of \(E\) such that
\[
L(E,s) = L(f,s).
\]
Equivalently, there exists a non-constant morphism
\[
X_0(N) \to E
\]
defined over \(\Q\).
\end{theorem}

An elliptic curve \(E\) over \(\Q\) has \textbf{semistable reduction} at a prime \(p\) if the reduction of \(E\) modulo \(p\) is either smooth (good reduction) or a nodal cubic (multiplicative reduction); in particular, additive reduction is not allowed. A curve is semistable if it has semistable reduction at all primes.

The Frey curve \(E_{a,b}\) constructed from a putative Fermat solution is semistable. Hence, by Wiles's theorem, it would be modular. But Ribet's theorem shows that a semistable Frey curve cannot be modular. Therefore, no such curve can exist, and Fermat's Last Theorem follows.

\begin{theorem}[Fermat's Last Theorem, Wiles 1994]
For any integer \(n \geq 3\), the Diophantine equation
\[
x^n + y^n = z^n
\]
has no solutions in positive integers \(x, y, z\).
\end{theorem}

\begin{proof}
It suffices to prove the theorem for \(n = 4\) and for \(n = p\) an odd prime, since any \(n \geq 3\) is divisible by 4 or by an odd prime. The case \(n = 4\) was proved by Fermat. For an odd prime \(p\), suppose there exist positive integers \(a, b, c\) with \(a^p + b^p = c^p\). Construct the Frey curve
\[
E_{a,b}: \quad y^2 = x(x - a^p)(x + b^p).
\]
This curve is semistable over \(\Q\). By Wiles's theorem, \(E_{a,b}\) is modular. By Ribet's theorem, the associated mod \(p\) Galois representation would be modular of level 2, which is impossible by a theorem of Mazur. This contradiction shows that no nontrivial solution exists.
\end{proof}

The full Modularity Theorem, proving modularity for all elliptic curves over \(\Q\) (not just semistable ones), was subsequently proved by Breuil, Conrad, Diamond, and Taylor in 2001, building on Wiles's methods.

\section{Connection to the Langlands Program}

Wiles's proof is intimately connected to the Langlands program, a sweeping vision of connections between number theory, representation theory, and geometry proposed by Robert Langlands in the 1960s. The Langlands program predicts that Galois representations should correspond to automorphic representations, and the Shimura--Taniyama--Weil conjecture is a special case of this correspondence.

For \(\GL_2\) over \(\Q\), the Langlands correspondence establishes a bijection between:
\begin{itemize}
    \item Irreducible, continuous, odd Galois representations \(\rho: \Gal(\overline{\Q}/\Q) \to \GL_2(\C)\) (and more generally, \(\ell\)-adic representations), and
    \item Cuspidal automorphic representations of \(\GL_2(\mathbb{A}_\Q)\) (which correspond to modular forms).
\end{itemize}

The modularity of elliptic curves can be seen as the statement that the Galois representation on the Tate module of an elliptic curve corresponds to an automorphic representation coming from a modular form. This is the \(p\)-adic (or \(\ell\)-adic) case of the Langlands correspondence.

Wiles's proof provided the first major evidence for the Langlands correspondence in the \(p\)-adic setting. The methods he developed --- deformation rings, modularity lifting, and the Taylor--Wiles patching argument --- have become fundamental tools in the Langlands program. They have been extended to prove:
\begin{itemize}
    \item The full Modularity Theorem (Breuil--Conrad--Diamond--Taylor, 2001).
    \item The Sato--Tate conjecture (Clozel--Harris--Shepherd-Barron--Taylor, 2008).
    \item Many cases of the Langlands correspondence for \(\GL_n\) over number fields (Scholze, Harris--Taylor, and others).
    \item Serre's modularity conjecture (Khare--Wintenberger, 2008).
\end{itemize}

Wiles shared the 1995/96 Wolf Prize with Robert Langlands, recognizing not only Wiles's proof of FLT but also the deep ways in which his work vindicated the Langlands philosophy. The prize citation noted that Wiles's work demonstrated the profound interconnectedness of number theory that Langlands had foreseen.

\section{Iwasawa Theory in Wiles's Proof}

Iwasawa theory, developed by Kenkichi Iwasawa in the 1950s and 1960s, studies the growth of ideal class groups in towers of number fields, particularly the cyclotomic \(\Z_p\)-extension. A central object is the \textbf{Iwasawa algebra} \(\Z_p[[\Gamma]]\), where \(\Gamma = \Gal(\Q_\infty/\Q) \cong \Z_p\).

Wiles had deep expertise in Iwasawa theory from his work with John Coates and Barry Mazur. The \textbf{Mazur--Wiles theorem} (1984) proved the main conjecture of Iwasawa theory for cyclotomic fields, which relates the characteristic ideal of a certain Galois group to a \(p\)-adic \(L\)-function.

In the proof of Fermat's Last Theorem, Iwasawa theory appears in several critical ways:

\begin{enumerate}
    \item \textbf{Control of Selmer groups:} Iwasawa theory provides bounds on the size of Selmer groups for elliptic curves in cyclotomic towers. These bounds are essential for controlling the deformation rings.
    \item \textbf{The Euler system:} The Kolyvagin--Flach approach uses an Euler system derived from Heegner points and Siegel units. Iwasawa-theoretic methods are used to relate this Euler system to the structure of the Selmer group.
    \item \textbf{The basic inequality:} Wiles used Iwasawa theory to prove a fundamental inequality between the size of the deformation ring and the size of the Hecke algebra, a crucial step in the numerical criterion.
\end{enumerate}

The Galois cohomology calculations at the heart of Wiles's proof rely heavily on the theory of Selmer groups and their relationship to \(p\)-adic \(L\)-functions, which are the natural domain of Iwasawa theory.

\section{The Kolyvagin--Flach Method}

The Kolyvagin--Flach method, developed by Victor Kolyvagin and Matthias Flach in the late 1980s and early 1990s, provides a technique for bounding Selmer groups of elliptic curves using Euler systems.

An \textbf{Euler system} is a collection of cohomology classes
\[
c_m \in H^1(\Q(\mu_m), T),
\]
for varying integers \(m\), that satisfy compatibility relations under the norm maps \(\Q(\mu_{mn}) \to \Q(\mu_m)\). The classes are typically constructed from special values of \(L\)-functions, Heegner points, or units in cyclotomic fields.

Kolyvagin used Euler systems of Heegner points to bound the Selmer group of an elliptic curve \(E\) over \(\Q\), obtaining remarkable results on the Birch and Swinnerton-Dyer conjecture. Flach extended these ideas to study the Selmer group attached to the symmetric square of a modular form.

Wiles attempted to use the Kolyvagin--Flach method to bound the Selmer groups arising in his deformation-theoretic approach. The original 1993 proof contained a gap precisely in this part: the Euler system constructed by Wiles did not satisfy all the necessary compatibility conditions to give the required bound.

The resolution of this gap came from two insights:
\begin{enumerate}
    \item The use of the Taylor--Wiles patching method to avoid the need for the full strength of the Euler system argument.
    \item A refined analysis of the local properties of the cohomology classes using the theory of Fontaine and Lafaille.
\end{enumerate}

In the final proof, the role of the Euler system was reduced to establishing a specific inequality --- the so-called \textbf{basic inequality} --- which could be verified through careful local calculations rather than through the global Euler system. This was the key simplification that allowed Wiles and Taylor to fill the gap.

\section{Impact and Legacy}

\subsection{Resolution of Fermat's Last Theorem}

The most immediate impact of Wiles's work was the resolution of the most famous open problem in mathematics. Fermat's Last Theorem had captured the imagination of mathematicians and the public for over three centuries. Its proof was front-page news worldwide and brought mathematics into the public spotlight in a way that few other achievements have.

\subsection{Transformation of Number Theory}

Wiles's proof revolutionized number theory. The techniques he developed --- deformation rings, modularity lifting, and the Taylor--Wiles patching method --- have become central tools in arithmetic geometry. They have been extended and generalized to prove:

\begin{itemize}
    \item The full Shimura--Taniyama--Weil conjecture (Breuil--Conrad--Diamond--Taylor, 2001).
    \item Serre's modularity conjecture (Khare--Wintenberger, 2008).
    \item The Sato--Tate conjecture (Clozel--Harris--Shepherd-Barron--Taylor, 2008).
    \item Potential automorphy theorems (Barnet-Lamb--Gee--Geraghty--Taylor, 2010+).
\end{itemize}

\subsection{The Langlands Program}

Wiles's proof provided the first major evidence that the Langlands correspondence for \(\GL_2\) could be proven using geometric and deformation-theoretic methods. The modularity lifting technique has been a central thread in the subsequent development of the Langlands program, leading to the proof of the fundamental lemma by Ngo Bao Chau and the development of the theory of perfectoid spaces by Peter Scholze.

\subsection{Recognition}

Andrew Wiles has received numerous awards for his work:

\begin{itemize}
    \item Wolf Prize in Mathematics (1995/96, shared with Robert Langlands).
    \item Schock Prize (1995).
    \item Ostrowski Prize (1995).
    \item Fermat Prize (1995).
    \item Royal Medal (1996).
    \item Wolf Prize (1996).
    \item MacArthur Fellowship (1997).
    \item Abel Prize (2016).
    \item Knight Commander of the Order of the British Empire (KBE, 2011).
\end{itemize}

He was elected a Fellow of the Royal Society (1989) and a Foreign Associate of the National Academy of Sciences (1997). He is currently the Regius Professor of Mathematics at Oxford University (since 2018) and a Fellow of Merton College, Oxford.

\subsection{The Full Modularity Theorem}

The complete proof of the Shimura--Taniyama--Weil conjecture (now the Modularity Theorem) by Breuil, Conrad, Diamond, and Taylor (2001) extended Wiles's methods to cover all elliptic curves over \(\Q\), not just semistable ones. This required handling curves with additive reduction at some primes, which necessitated a more sophisticated deformation-theoretic analysis.

The Modularity Theorem states:

\begin{theorem}[Modularity Theorem, 2001]
Every elliptic curve over \(\Q\) is modular. For any elliptic curve \(E/\Q\) of conductor \(N\), there exists a newform \(f\) of weight 2 and level \(N\) such that \(L(E,s) = L(f,s)\).
\end{theorem}

This theorem is one of the most important results in 21st-century mathematics. It has far-reaching consequences for the arithmetic of elliptic curves, including the Birch and Swinnerton-Dyer conjecture, the congruent number problem, and the construction of rational points on elliptic curves.

\section{Timeline}

\begin{tabularx}{\linewidth}{|l|X|}
\hline
\textbf{Year} & \textbf{Event} \\ \hline
1953 & Born on 11 April in Cambridge, England \\ \hline
1971 & Entered Merton College, Oxford University \\ \hline
1974 & First-class degree in Mathematics, Oxford University \\ \hline
1977 & Junior Research Fellow, Clare College, Cambridge \\ \hline
1979 & Ph.D. at Clare College, Cambridge (advisor: John Coates) \\ \hline
1980 & Visiting Scholar, Institute for Advanced Study, Princeton \\ \hline
1981 & Assistant Professor, Harvard University \\ \hline
1982 & Professor, Princeton University \\ \hline
1986 & Ribet proves that FLT follows from modularity of semistable elliptic curves; Wiles begins secret work on the proof \\ \hline
1989 & Elected Fellow of the Royal Society \\ \hline
1993 & Announces proof of Fermat's Last Theorem at Isaac Newton Institute, Cambridge (June 21--23); gap discovered in the proof \\ \hline
1994 & With Richard Taylor, fixes the gap (September--October); submits final papers to the Annals of Mathematics \\ \hline
1995 & Papers published in the Annals of Mathematics; proof accepted as correct \\ \hline
1995/96 & Awarded the Wolf Prize in Mathematics (shared with Robert Langlands) \\ \hline
1997 & Awarded the MacArthur Fellowship \\ \hline
2001 & Breuil--Conrad--Diamond--Taylor prove full Modularity Theorem \\ \hline
2011 & Appointed Knight Commander of the Order of the British Empire (KBE) \\ \hline
2016 & Awarded the Abel Prize \\ \hline
2018 & Appointed Regius Professor of Mathematics at Oxford University \\ \hline
\end{tabularx}

\section{Frequently Asked Questions}

\subsection*{Q1: Did Fermat actually have a proof of his Last Theorem?}

Most historians and mathematicians believe that Fermat did not have a correct proof. Fermat's claimed proof, if it existed, would have been based on the method of infinite descent, which works for \(n = 4\) but cannot be extended to all exponents. The difficulty of the problem and the 350 years it took to solve it strongly suggest that Fermat was mistaken in believing he had a proof.

\subsection*{Q2: How long did Wiles work on the proof?}

Wiles worked in secrecy on Fermat's Last Theorem for approximately seven years, from 1986 (when Ribet proved the connection between FLT and modularity) to 1993 (when he announced the proof). After the gap was discovered, he worked for another year (with Richard Taylor) to fix it. In total, the proof took about eight years of intensive work.

\subsection*{Q3: What was the gap in Wiles's 1993 proof?}

The gap was in the portion of the proof concerning the use of the Kolyvagin--Flach method to bound Selmer groups. Specifically, the construction of an Euler system used to prove the basic inequality was incomplete. The gap was fixed by Wiles and Taylor using an extended version of the Taylor--Wiles patching argument, which avoided the need for the problematic Euler system.

\subsection*{Q4: What is the relationship between Wiles's proof and the Langlands program?}

Wiles's proof is a spectacular confirmation of the Langlands philosophy. The modularity of elliptic curves (the Shimura--Taniyama--Weil conjecture) is a special case of the Langlands correspondence for \(\GL_2\). Wiles's methods --- particularly modularity lifting and deformation rings --- have become fundamental tools in the Langlands program, leading to the proof of Serre's modularity conjecture, the Sato--Tate conjecture, and many other results.

\subsection*{Q5: What is the current status of the Shimura--Taniyama--Weil conjecture?}

The full conjecture, now called the Modularity Theorem, was proved in 2001 by Breuil, Conrad, Diamond, and Taylor, building on Wiles's methods. The theorem states that every elliptic curve over \(\Q\) is modular. This is a central result in modern number theory.

\subsection*{Q6: What is an elliptic curve, and why is it relevant to FLT?}

An elliptic curve is a smooth projective curve of genus 1 with a distinguished point, typically given by an equation \(y^2 = x^3 + ax + b\). Elliptic curves are relevant to FLT because Gerhard Frey showed that a solution to \(a^p + b^p = c^p\) would give an elliptic curve (the Frey curve) with extraordinary properties. Ribet proved that such a curve cannot be modular, and Wiles proved that every semistable elliptic curve is modular, giving a contradiction unless no Fermat solution exists.

\subsection*{Q7: What is modularity?}

Modularity is the property that an elliptic curve corresponds to a modular form, a highly symmetric holomorphic function on the upper half-plane. Modular forms have deep connections to number theory, and the correspondence allows problems about elliptic curves (like counting rational points) to be translated into problems about modular forms (like analyzing Fourier coefficients).

\subsection*{Q8: What are the major open problems related to Wiles's work?}

Several deep problems remain:
\begin{enumerate}
    \item \textbf{The Birch and Swinnerton-Dyer conjecture:} This conjecture relates the rank of an elliptic curve to the behavior of its \(L\)-function. Wiles's methods have been used to make progress on this conjecture, but a full proof remains elusive.
    \item \textbf{The Langlands correspondence for \(\GL_n\) over number fields:} While much progress has been made, the full functoriality conjectures of Langlands remain open.
    \item \textbf{The modularity of abelian surfaces over \(\Q\):} Extending modularity to higher-dimensional abelian varieties is a major open problem.
    \item \textbf{Generalizations of the Taylor--Wiles method to other contexts:} While the method has been successfully extended to many settings, further generalizations are needed to prove the full Langlands correspondence.
\end{enumerate}

\subsection*{Q9: How did Wiles's upbringing influence his mathematical career?}

Wiles grew up in Cambridge, the son of a theology professor at Oxford. He has said that he was drawn to Fermat's Last Theorem at age ten when he encountered the problem in a library book. This early fascination set him on a path toward number theory, and his PhD work with John Coates on Iwasawa theory and elliptic curves provided the technical foundation for his later proof.

\subsection*{Q10: Why is the proof of FLT considered historically important beyond mathematics?}

The proof of Fermat's Last Theorem is considered a landmark of human intellectual achievement because it solved a problem that had challenged the greatest minds for over 350 years. It represents the culmination of centuries of mathematical development and the triumph of a deep, structural vision over a seemingly elementary question. It also demonstrated the profound interconnectedness of mathematics, bringing together disparate fields in unexpected ways.

\section{Related Chapters}

See also Chapter~04 (Weil) on number theory and algebraic geometry.

\section*{Exercises}
\addcontentsline{toc}{section}{Exercises}

\begin{enumerate}
    \item [B] Show that if Fermat's Last Theorem holds for exponent \(n\), it holds for any multiple of \(n\). Conclude that it suffices to prove FLT for \(n = 4\) and for odd primes.
    \item [I] Write the Frey curve \(E_{a,b}\) in the standard Weierstrass form and compute its discriminant and \(j\)-invariant. Show that the curve is semistable.
    \item [I] Prove that the mod \(p\) Galois representation \(\overline{\rho}_{E,p}\) of the Frey curve is irreducible for \(p \geq 5\).
    \item [I] Show that the conductor of the Frey curve \(E_{a,b}\) is \(\prod_{\ell \mid abc} \ell\).
    \item [A] Verify Ribet's level-lowering theorem for the specific case where the Frey curve is associated to a Fermat solution. Show that the level can be reduced to 2.
    \item [B] Prove that there are no non-zero cusp forms of weight 2 and level 2. (Hint: use the dimension formula for spaces of modular forms.)
    \item [I] Let \(E\) be an elliptic curve over \(\Q\) with \(j\)-invariant \(j\). Show that if \(j \notin \Z\), then \(E\) has potential good reduction at some prime.
    \item [I] Show that the deformation ring \(R\) for a mod \(p\) representation of \(\Gal(\overline{\Q}/\Q)\) is a complete intersection if the representation is unobstructed.
    \item [I] State the main conjecture of Iwasawa theory for cyclotomic fields. Explain its relationship to the Mazur--Wiles theorem.
    \item [I] Let \(\overline{\rho}: \Gal(\overline{\Q}/\Q) \to \GL_2(\F_p)\) be a continuous representation. Define the universal deformation ring \(R\) and show that it is a complete Noetherian local \(\Z_p\)-algebra.
    \item [I] Show that the Hecke algebra \(\mathbb{T}\) generated by Hecke operators acting on \(S_2(\Gamma_0(N))\) is a reduced \(\Z\)-algebra.
    \item [B] Prove that if a prime \(\ell\) divides the conductor of an elliptic curve \(E/\Q\), then the curve has bad reduction at \(\ell\).
    \item [I] Explain why the Taylor--Wiles patching argument requires the introduction of auxiliary primes \(\ell_i \equiv 1 \pmod{p^N}\).
    \item [B] Let \(f\) be a newform of weight 2 and level \(N\). Show that the Fourier coefficients \(a_p(f)\) satisfy \(|a_p(f)| \leq 2\sqrt{p}\) (the Hasse bound).
    \item [I] Verify that if \(a^p + b^p = c^p\), then the Frey curve has discriminant \(\Delta = (abc)^{2p}/256\) and compute its minimal discriminant.
    \item [I] Show that the Fontaine--Mazur conjecture predicts that Galois representations arising from geometry should be potentially semistable at all primes.
    \item [I] Prove that the set of rational points on an elliptic curve forms a finitely generated abelian group (Mordell's theorem). State the weak Mordell--Weil theorem and explain its role.
    \item [I] Let \(E/\Q\) be an elliptic curve with complex multiplication. Show that \(E\) is modular. (This was proved by Shimura before Wiles's work.)
    \item [I] Consider the elliptic curve \(E: y^2 + y = x^3 - x^2\). Show that \(E\) is modular by finding the corresponding modular form.
\end{enumerate}

\section*{Glossary}
\addcontentsline{toc}{section}{Glossary}

\begin{description}
    \item[Abelian variety] A projective algebraic variety that is also an algebraic group; in dimension 1, an elliptic curve.
    \item[Automorphic form] A generalization of modular forms to higher-dimensional groups, central to the Langlands program.
    \item[Complete intersection] A local ring that can be expressed as a quotient of a regular local ring by a regular sequence.
    \item[Conductor] An integer associated to an elliptic curve that encodes its bad reduction behavior.
    \item[Cusp form] A modular form that vanishes at all cusps of the modular curve.
    \item[Deformation ring] A universal object parameterizing lifts of a residual Galois representation.
    \item[Elliptic curve] A smooth projective curve of genus 1 with a distinguished rational point.
    \item[Euler system] A compatible system of cohomology classes used to bound Selmer groups.
    \item[Frey curve] The elliptic curve \(y^2 = x(x-a^p)(x+b^p)\) associated to a putative solution of Fermat's equation.
    \item[Galois representation] A continuous homomorphism from \(\Gal(\overline{\Q}/\Q)\) to \(\GL_n(R)\) for some ring \(R\).
    \item[Hecke algebra] The algebra generated by Hecke operators acting on a space of modular forms.
    \item[Hecke operator] A linear operator on spaces of modular forms defined by double coset decompositions, fundamental for the theory.
    \item[Iwasawa algebra] The completed group ring \(\Z_p[[\Gamma]]\) where \(\Gamma \cong \Z_p\).
    \item[Iwasawa theory] The study of class groups in cyclotomic \(\Z_p\)-extensions and their relationship to \(p\)-adic \(L\)-functions.
    \item[\(j\)-invariant] A modular function that classifies elliptic curves up to isomorphism over \(\C\).
    \item[Langlands program] A vast network of conjectures relating Galois representations to automorphic forms.
    \item[Level (of a modular form)] The integer \(N\) such that the modular form is invariant under \(\Gamma_0(N)\).
    \item[Modular form] A holomorphic function on the upper half-plane satisfying a functional equation under the action of a congruence subgroup.
    \item[Modular curve] The quotient of the upper half-plane by a congruence subgroup, an algebraic curve defined over \(\Q\).
    \item[Modularity] The property of an elliptic curve being associated to a modular form via the Eichler--Shimura construction.
    \item[Modularity lifting] The principle that if a mod \(p\) Galois representation is modular, then its \(p\)-adic lift is also modular.
    \item[Newform] A cusp form that is a normalized eigenform for all Hecke operators and is not induced from a lower level.
    \item[Residual representation] The reduction of a \(p\)-adic Galois representation modulo \(p\).
    \item[Sato--Tate conjecture] The conjecture that the distribution of the normalized Fourier coefficients \(a_p/(2\sqrt{p})\) follows the Sato--Tate measure.
    \item[Selmer group] A Galois cohomology group classifying extensions of Galois modules with prescribed local behavior.
    \item[Semistable reduction] Reduction of an elliptic curve that is either smooth or nodal (multiplicative) at a prime.
    \item[Taylor--Wiles method] The patching technique used to prove that deformation rings and Hecke algebras are isomorphic.
    \item[Weierstrass equation] A cubic equation of the form \(y^2 = x^3 + ax + b\) defining an elliptic curve.
\end{description}
